\section{The \texorpdfstring{$[[5,1,2]]_3$}{[[5,1,2]]3} Code}
\label{sub:code}\label{sec: error detection code}\label{sec:method}

% Sec. II of the restructured sim version (2026-11-19, user: simulation paper around the code design, LRB-D and the QPU layouts; style of the SNAP draft, jpw-review/sim/style-snap.md). Built from the former Secs. 2.1 and 2.2 ("method - conventions.tex", "method - [[5,1,2]]_3 qutrit surface code.tex", kept unused for reference). Outline: jpw-review/sim/outline-II-IV.md.
Let us begin with what a code must provide for a benchmark of logical gates on a small device. It must use few qutrits, so that the encoded circuits fit on the device. It must detect every error on a single qutrit, so that postselection can reject the shots, the executions of a circuit, in which such an error occurs. And its logical Clifford gates must cost few physical gates, so that a logical gate is not much longer than its physical counterpart. The $[[5,1,2]]_3$ code, a patch of the qutrit surface code~\cite{Moussa2016}, meets all three requirements, and this section describes it.

\label{sub:conventions}We first fix the notation. We work with qudits of prime local dimension $q$, with $q = 3$ (qutrits) except in the comparison with qubits, where $q = 2$. Let $\omega = e^{2\pi i/q}$, and let $\ket{j}$ with $j \in \{0, \dots, q-1\}$ be the computational basis. The Pauli operators $X$ and $Z$ and the Clifford gates $H$, $S$, SUM and CZ act as
\begin{equation}
  \begin{gathered}
    \begin{aligned}
      X\ket{j} &= \ket{j+1}, & Z\ket{j} &= \omega^{j}\ket{j}, \\
      S\ket{j} &= \omega^{j(j-1)/2}\ket{j}, & H\ket{j} &= \frac{1}{\sqrt{q}}\sum_{l=0}^{q-1}\omega^{jl}\ket{l},
    \end{aligned} \\
    \begin{aligned}
      \mathrm{SUM}\,\ket{j}\ket{l} &= \ket{j}\ket{j+l}, \\
      \mathrm{CZ}\,\ket{j}\ket{l} &= \omega^{jl}\ket{j}\ket{l},
    \end{aligned}
  \end{gathered}
  \label{eq:conventions}
\end{equation}
with all arithmetic modulo $q$ (\appref{app:background} lists their conjugation rules). The form of $S$ holds for odd $q$. SUM adds the value of its first qutrit, the control, to its second, the target, and SWAP exchanges two qutrits. A subscript on a gate names the qutrit, or the pair of qutrits, it acts on. We measure in the $X$ basis by applying $H^{-1}$ and then measuring in the computational basis, so the outcome $j$ belongs to the $X$ eigenvalue $\omega^{-j}$.

\subsection{Stabilizers, Logical Operators and Distance}

The code stores one logical qutrit in five data qutrits $q_0, \dots, q_4$, which meets the first requirement. Its code space is the joint $+1$ eigenspace of four stabilizer generators of weight three, two of $Z$ type and two of $X$ type,
\begin{align}
  G_1 &= Z_0 Z_1^{-1} Z_3^{-1}, & G_2 &= Z_1 Z_2 Z_4^{-1}, \label{eq:s1}\\
  G_3 &= X_0 X_1 X_2^{-1}, & G_4 &= X_1^{-1} X_3 X_4^{-1} .
  \label{eq:s4}
\end{align}
We call the exponent of $\omega$ in the measured eigenvalue of $G_j$ the syndrome of $G_j$, and the syndromes are trivial when all four are zero, that is, when every eigenvalue is $+1$. The logical Pauli operators of the code are
\begin{equation}
  \bar Z = Z_0 Z_2, \qquad
  \bar X = X_0 X_3 .
  \label{eq:logicalPauli}
\end{equation}
Every one-qutrit Pauli error fails to commute with at least one generator and so gives a nontrivial syndrome, which meets the second requirement. The logical operators of Eq.~\eqref{eq:logicalPauli} have weight two, however, so two errors can together act as a logical operator without changing the syndrome. The distance is therefore $d = 2$, and a nontrivial syndrome reveals an error without identifying it.

\subsection{Logical Clifford Gates}

The third requirement is met because the logical Clifford gates are transversal in the sense of the folded surface code~\cite{Moussa2016}. They consist of one-qutrit gates and of two-qutrit gates on pairs of qutrits that the fold of this construction places on top of each other (\appref{app:background-folded}). The single-qutrit logical Clifford group is generated by
\begin{align}
  \bar H &=
  \left( H_0 H_1^{-1} H_2 H_3 H_4 \right) \mathrm{SWAP}_{2,3}, \label{eq:logical_H}\\
  \bar S &= S_0\, S_1^{-1}\, \mathrm{CZ}_{2,3}\, S_4 .
  \label{eq:logical_HS}
\end{align}
These expressions have a simple reading. Each generator applies one layer of one-qutrit gates, with the inverse gate on $q_1$, and a single two-qutrit gate on the pair $q_2$, $q_3$: a SWAP in $\bar H$ and a CZ in $\bar S$. Both generators map the stabilizer group to itself and act on $\bar X$ and $\bar Z$ as $H$ and $S$ act on $X$ and $Z$ (\appref{app:rounds}). The words of the logical Cliffords that we simulate also contain the logical Paulis (\appref{sec:clifford-gen}), which are products of one-qutrit gates, so the logical Cliffords use a single coupling between data qutrits, the one between $q_2$ and $q_3$. This coupling is the first thing that a layout of the code on a device must provide (\autoref{sub:ancilla_layouts}).


% ============================================================
% fig - 5-qutrit code.tex
% Figure: [[5,1,2]]_3 patch definitions, logical Cliffords, stabilizers,
% and logical Paulis (d=2 nearest-neighbour construction).
% All macros come from 'tikz preamble.tex'.
% ============================================================

\begin{figure*}[tp]
\centering
% ===================== (a) Reference =====================
\begin{subfigure}[b]{0.9\textwidth}
\centering
\begin{tikzpicture}[line cap=round,line join=round,thick,scale=1.0]
  \coordinate (Left)   at (-4.85, 0);
  \coordinate (Middle) at (-1.85, 0);
  \coordinate (XGrid)  at ( 1.55, 0);
  \coordinate (ZGrid)  at ( 4.45, 0);

  \node[above,font=\small] at ($(Left)+(0,1.72)$) {Unfolded};
  \node[above,font=\small] at ($(Middle)+(0,1.72)$) {Folded};
  \node[above,font=\small] at ($(XGrid)!0.5!(ZGrid)+(0,1.80)$) {Redefined Stabilizers};
  \node[above,font=\footnotesize,text=Xcol!80!black] at ($(XGrid)+(0,1.20)$) {$X$-type Stabilizers};
  \node[above,font=\footnotesize,text=Zcol!85!black] at ($(ZGrid)+(0,1.20)$) {$Z$-type Stabilizers};
  \node[font=\scriptsize,text=Xcol!80!black] at ($(XGrid)+(0,-1.25)$) {$G_3,\;G_4$};
  \node[font=\scriptsize,text=Zcol!85!black] at ($(ZGrid)+(0,-1.25)$) {$G_1,\;G_2$};

  % --- Left Patch (Unfolded) ---
  \begin{scope}[shift={(Left)}]
    \FiveQPatchFilled
    \FiveQPatchCircles
  \end{scope}

  % --- Middle Patch (Folded) ---
  \begin{scope}[shift={(Middle)}]
    \FiveQPatchOverlapped
    \def\m{0.30}
    \draw[dashed, draw=red, line width=1pt, dash pattern=on 3pt off 2pt]
      ($(q0)+(-\m,\m)$) -- ($(q4)+(\m,-\m)$);
    \FiveQPatchCirclesOneUnderThree
  \end{scope}

  \draw[-{Stealth[length=3mm,width=2mm]},red,line width=1.6pt]
    ($(Middle)+(1.35,0)$) -- ($(XGrid)+(-1.35,0)$);

  % --- Redefined X-type grid ---
  \begin{scope}[shift={(XGrid)}]
    \FiveQPatchRedefinedXGrid
    \FiveQPatchCircles
  \end{scope}

  % --- Redefined Z-type grid ---
  \begin{scope}[shift={(ZGrid)}]
    \FiveQPatchRedefinedZGrid
    \FiveQPatchCircles
  \end{scope}
\end{tikzpicture}
\caption{Five-qutrit patch}
\end{subfigure}

\vspace{2.0ex}

% ===================== (b) Definition =====================
\hspace*{-1.5em}
\begin{subfigure}[c]{0.48\textwidth}
\centering
\begin{tikzpicture}[line cap=round,line join=round,thick,scale=1.0]
  \coordinate (LeftTop) at (-2.2,  1.50);
  \coordinate (LeftBot) at (-2.2, -1.50);
  \coordinate (RightTop) at ( 2.2,  1.50);
  \coordinate (RightBot) at ( 2.2, -1.50);

  \node[scale=1.0, anchor=east] at ($(LeftTop)+(-1.4, -1.5)$) {$\bar{H'}\!:=$};
  \node[scale=1.0, anchor=east] (HLabel) at ($(RightTop)+(-1.3, -1.5)$) {$\bar{H}\!:=$};

  \def\lw{1.5pt} \def\rc{8pt}
  % LEFT BOX 
  \path let \p1 = ($(LeftTop)+(-1.3,1.7)$), \p2 = ($(LeftBot)+(1.4,-1.4)$)
        in coordinate (LTL) at (\x1,\y1) coordinate (LBR) at (\x2,\y2)
           coordinate (LTR) at (\x2,\y1) coordinate (LBL) at (\x1,\y2);
  \coordinate (LMidLeft) at ($(LTL)!0.5!(LBL)$); \coordinate (LMidRight) at ($(LTR)!0.5!(LBR)$);
  \draw[HCol, dash pattern=on 2pt off 5pt, line width=\lw, rounded corners=\rc] (LMidLeft)--(LTL)--(LTR)--(LMidRight);
  \draw[SwapCol, dash pattern=on 2pt off 5pt, line width=\lw, rounded corners=\rc] (LMidRight)--(LBR)--(LBL)--(LMidLeft);

  % RIGHT BOX
  \path let \p1 = ($(RightTop)+(-1.3,1.7)$), \p2 = ($(RightBot)+(1.4,-1.4)$)
        in coordinate (RTL) at (\x1,\y1) coordinate (RBR) at (\x2,\y2)
           coordinate (RTR) at (\x2,\y1) coordinate (RBL) at (\x1,\y2);
  \coordinate (RMidLeft) at ($(RTL)!0.5!(RBL)$); \coordinate (RMidRight) at ($(RTR)!0.5!(RBR)$);
  \draw[HCol, dash pattern=on 2pt off 5pt, line width=\lw, rounded corners=\rc] (RMidLeft)--(RTL)--(RTR)--(RMidRight);
  \draw[SwapCol, dash pattern=on 2pt off 5pt, line width=\lw, rounded corners=\rc] (RMidRight)--(RBR)--(RBL)--(RMidLeft);

  \begin{scope}[shift={(LeftTop)}]
    \FiveQPatchEdges
    \foreach \P/\txt in {q0/0,q1/1,q2/2,q3/3,q4/4}{ \draw[line width=\qdLine,fill=HCol] (\P) circle (\qdR); \node[font=\numFont, text=white] at (\P) {\txt}; }
    \WedgeArrow{q0}{0}\WedgeArrow{q1}{0}\WedgeArrow{q3}{0}\WedgeArrow{q4}{0}\WedgeArrow{q2}{90}
    \node[anchor=south, font=\small, yshift=0mm] at (0, 1.0) {Layer 1};
    \begin{scope}[text=HCol, font=\footnotesize]
       \node[below, yshift=-3pt] at (q0) {$H_0$}; \node[below, yshift=-3pt] at (q1) {$H_1$};
       \node[right, xshift=2pt]  at (q2) {$H_2^{-1}$}; \node[above, yshift=3pt]  at (q3) {$H_3$}; \node[above, yshift=3pt]  at (q4) {$H_4$};
    \end{scope}
  \end{scope}

  \begin{scope}[shift={(LeftBot)}]
    \FiveQPatchEdges
    \foreach \P/\txt in {q0/0,q2/2,q4/4}{ \draw[line width=\qdLine,fill=white] (\P) circle (\qdR); \node[font=\numFont] at (\P) {\txt}; }
    \foreach \P/\txt in {q1/1,q3/3}{ \draw[line width=\qdLine,fill=SwapCol] (\P) circle (\qdR); \node[font=\numFont, text=white] at (\P) {\txt}; }
    \WedgeArrow{q0}{0}\WedgeArrow{q1}{0}\WedgeArrow{q3}{0}\WedgeArrow{q4}{0}\WedgeArrow{q2}{90}
    \node[anchor=south, font=\small, yshift=0mm] at (0, 1.0) {Layer 2};
    \draw[{Stealth[length=3mm,width=2mm]}-{Stealth[length=3mm,width=2mm]}, color=SwapCol, line width=1.2pt, shorten >=5pt, shorten <=3pt]
      (q1) to[bend left=25] node[midway, below right, rotate=52, xshift=-7pt, font=\small] {SWAP} (q3);
  \end{scope}

  \draw[-{Stealth[length=3mm,width=2mm]},red!75,line width=1.6pt] (-0.6, 0) -- (-0.1, 0);

  \begin{scope}[shift={(RightTop)}]
    \FiveQPatchEdges
    \foreach \P/\txt in {q0/0,q1/1,q2/2,q3/3,q4/4}{ \draw[line width=\qdLine,fill=HCol] (\P) circle (\qdR); \node[font=\numFont, text=white] at (\P) {\txt}; }
    \WedgeArrow{q0}{0}\WedgeArrow{q1}{0}\WedgeArrow{q3}{0}\WedgeArrow{q4}{0}\WedgeArrow{q2}{90}
    \node[anchor=south, font=\small, yshift=0mm] at (0, 1.0) {Layer 1};
    \begin{scope}[text=HCol, font=\footnotesize]
       \node[below, yshift=-3pt] at (q0) {$H_0$}; \node[below, yshift=-3pt] at (q1) {$H_1^{-1}$};
       \node[right, xshift=2pt]  at (q2) {$H_2$}; \node[above, yshift=3pt]  at (q3) {$H_3$}; \node[above, yshift=3pt]  at (q4) {$H_4$};
    \end{scope}
  \end{scope}

  \begin{scope}[shift={(RightBot)}]
    \FiveQPatchEdges
    \foreach \P/\txt in {q0/0,q1/1,q4/4}{ \draw[line width=\qdLine,fill=white] (\P) circle (\qdR); \node[font=\numFont] at (\P) {\txt}; }
    \foreach \P/\txt in {q2/2,q3/3}{ \draw[line width=\qdLine,fill=SwapCol] (\P) circle (\qdR); \node[font=\numFont, text=white] at (\P) {\txt}; }
    \WedgeArrow{q0}{0}\WedgeArrow{q1}{0}\WedgeArrow{q3}{0}\WedgeArrow{q4}{0}\WedgeArrow{q2}{90}
    \node[anchor=south, font=\small, yshift=0mm] at (0, 1.0) {Layer 2};
    \draw[{Stealth[length=3mm,width=2mm]}-{Stealth[length=3mm,width=2mm]}, color=SwapCol, line width=1.2pt, shorten >=3pt, shorten <=3pt]
      (q2) to[bend right=45] node[midway, below, xshift =-7pt, yshift=14pt, rotate=41, font=\small] {SWAP} (q3);
  \end{scope}
\end{tikzpicture}
\caption{Logical $\bar H$}
\label{fig: logical H}
\end{subfigure}
\quad
% ===================== (c) S Definition =====================
\begin{subfigure}[c]{0.48\textwidth}
\centering
\begin{tikzpicture}[line cap=round,line join=round,thick,scale=1.0]
  \coordinate (Left)  at (-2.2, 0);
  \coordinate (Right) at ( 2.2, 0);
  \def\lw{1.5pt} \def\rc{8pt}

  \coordinate (LTL) at ($(Left)+(-1.3, 1.4)$); \coordinate (LBR) at ($(Left)+(1.4, -1.4)$);
  \coordinate (LTR) at ($(Left)+(1.4, 1.4)$);  \coordinate (LBL) at ($(Left)+(-1.3, -1.4)$);
  \coordinate (LMidLeft) at ($(LTL)!0.5!(LBL)$); \coordinate (LMidRight) at ($(LTR)!0.5!(LBR)$);
  \draw[HCol, dash pattern=on 2pt off 5pt, line width=\lw, rounded corners=\rc] (LMidLeft)--(LTL)--(LTR)--(LMidRight);
  \draw[SwapCol, dash pattern=on 2pt off 5pt, line width=\lw, rounded corners=\rc] (LMidRight)--(LBR)--(LBL)--(LMidLeft);

  \coordinate (RTL) at ($(Right)+(-1.3, 1.4)$); \coordinate (RBR) at ($(Right)+(1.4, -1.4)$);
  \coordinate (RTR) at ($(Right)+(1.4, 1.4)$);  \coordinate (RBL) at ($(Right)+(-1.3, -1.4)$);
  \coordinate (RMidLeft) at ($(RTL)!0.5!(RBL)$); \coordinate (RMidRight) at ($(RTR)!0.5!(RBR)$);
  \draw[HCol, dash pattern=on 2pt off 5pt, line width=\lw, rounded corners=\rc] (RMidLeft)--(RTL)--(RTR)--(RMidRight);
  \draw[SwapCol, dash pattern=on 2pt off 5pt, line width=\lw, rounded corners=\rc] (RMidRight)--(RBR)--(RBL)--(RMidLeft);

  \node[scale=1.0, anchor=east] at ($(Left)+(-1.4, 0)$) {$\bar{S'}\!:=$};
  \node[scale=1.0, anchor=east] (SLabel) at ($(Right)+(-1.4, 0)$) {$\bar{S}\!:=$};

  \newcommand{\DrawLHSPattern}{
      \FiveQPatchEdges
      \foreach \P in {q0, q2, q4} { \draw[line width=\qdLine,fill=HCol] (\P) circle (\qdR); }
      \foreach \P in {q1, q3} { \draw[line width=\qdLine,fill=SwapCol] (\P) circle (\qdR); }
      \foreach \P/\txt in {q0/0,q1/1,q2/2,q3/3,q4/4}{ \node[font=\numFont, text=white] at (\P) {\txt}; }
      \WedgeArrow{q0}{0}\WedgeArrow{q1}{0}\WedgeArrow{q3}{0}\WedgeArrow{q4}{0}\WedgeArrow{q2}{90}
      \node[text=HCol, font=\small, below, yshift=-3pt] at (q0) {$S_0$};
      \node[text=HCol, font=\small, right, xshift=3pt] at (q2) {$S_2^{-1}$};
      \node[text=HCol, font=\small, above, yshift=3pt] at (q4) {$S_4$};
      \draw[color=SwapCol, line width=1.2pt, shorten <=4pt, shorten >=4pt] 
         (q1) to[out=180, in=90] coordinate[pos=0.45] (mid) (q3);
      \node[text=SwapCol, font=\small, rotate=52, anchor=south, xshift=-12pt, yshift=-1pt] at (mid) {$CZ_{1\to 3}$};
  }
  \newcommand{\DrawRHSPattern}{
      \FiveQPatchEdges
      \foreach \P in {q0, q1, q4} { \draw[line width=\qdLine,fill=HCol] (\P) circle (\qdR); }
      \foreach \P in {q2, q3} { \draw[line width=\qdLine,fill=SwapCol] (\P) circle (\qdR); }
      \foreach \P/\txt in {q0/0,q1/1,q2/2,q3/3,q4/4}{ \node[font=\numFont, text=white] at (\P) {\txt}; }
      \WedgeArrow{q0}{0}\WedgeArrow{q1}{0}\WedgeArrow{q3}{0}\WedgeArrow{q4}{0}\WedgeArrow{q2}{90}
      \node[text=HCol, font=\small, below, yshift=-3pt] at (q0) {$S_0$};
      \node[text=HCol, font=\small, below, yshift=-3pt] at (q1) {$S_1^{-1}$};
      \node[text=HCol, font=\small, above, yshift=3pt] at (q4) {$S_4$};
      \draw[color=SwapCol, line width=1.2pt, shorten <=4pt, shorten >=4pt] 
         (q3) to[out=90, in=180] coordinate[pos=0.5] (midR) (q2);
      \node[text=SwapCol, font=\small, rotate=40, anchor=south, xshift=5pt, yshift=-2pt] at (midR) {$CZ_{2\to 3}$};
  }
  \begin{scope}[shift={(Left)}] \DrawLHSPattern \end{scope}
  \begin{scope}[shift={(Right)}] \DrawRHSPattern \end{scope}
  \draw[-{Stealth[length=3mm,width=2mm]},red,line width=1.6pt] (-0.6, 0) -- (-0.1, 0);
\end{tikzpicture}
\caption{Logical $\bar S$}
\label{fig: logical S}
\end{subfigure}

\vspace{2.0ex}

% ===================== (d) Parallel Operations =====================
\begin{subfigure}[b]{0.9\textwidth}
\centering
\begin{tikzpicture}[line cap=round,line join=round,thick,scale=1.0]
  \def\yTop{1.6} \def\yBot{-1.2}

  % COL 1
  \def\x{-4.8}
  \begin{scope}[shift={(\x, \yTop)}]
     \FiveQPatchEdges
     \foreach \P/\txt in {q1/1, q4/4}{ \draw[line width=\qdLine,fill=white] (\P) circle (\qdR); \node[font=\numFont] at (\P) {\txt}; }
     \foreach \P/\txt in {q0/0, q2/2, q3/3}{ \draw[line width=\qdLine,fill=Zcol] (\P) circle (\qdR); \node[font=\numFont, text=white] at (\P) {\txt}; }
     \WedgeArrow{q0}{0}\WedgeArrow{q1}{0}\WedgeArrow{q3}{0}\WedgeArrow{q4}{0}\WedgeArrow{q2}{90}
     \begin{scope}[text=Zcol!85!black, font=\footnotesize]
        \node[below, yshift=-3pt] at (q0) {$Z_0$}; \node[right, xshift=2pt] at (q2) {$Z_2^{-1}$}; \node[above, yshift=3pt] at (q3) {$Z_3^{-1}$};
     \end{scope}
  \end{scope}
  \begin{scope}[shift={(\x, \yBot)}]
     \FiveQPatchEdges
     \foreach \P/\txt in {q2/2, q4/4}{ \draw[line width=\qdLine,fill=white] (\P) circle (\qdR); \node[font=\numFont] at (\P) {\txt}; }
     \foreach \P/\txt in {q0/0, q1/1, q3/3}{ \draw[line width=\qdLine,fill=Zcol] (\P) circle (\qdR); \node[font=\numFont, text=white] at (\P) {\txt}; }
     \WedgeArrow{q0}{0}\WedgeArrow{q1}{0}\WedgeArrow{q3}{0}\WedgeArrow{q4}{0}\WedgeArrow{q2}{90}
     \begin{scope}[text=Zcol!85!black, font=\footnotesize]
        \node[below, yshift=-3pt] at (q0) {$Z_0$}; \node[below, yshift=-3pt] at (q1) {$Z_1^{-1}$}; \node[above, yshift=3pt] at (q3) {$Z_3^{-1}$};
     \end{scope}
  \end{scope}
  \draw[-{Stealth[length=3mm,width=2mm]},red!75,line width=1.6pt] (\x, 0.425) -- (\x, -0.1);

  % COL 2
  \def\x{-1.6}
  \begin{scope}[shift={(\x, \yTop)}]
     \FiveQPatchEdges
     \foreach \P/\txt in {q0/0, q3/3}{ \draw[line width=\qdLine,fill=white] (\P) circle (\qdR); \node[font=\numFont] at (\P) {\txt}; }
     \foreach \P/\txt in {q1/1, q2/2, q4/4}{ \draw[line width=\qdLine,fill=Zcol] (\P) circle (\qdR); \node[font=\numFont, text=white] at (\P) {\txt}; }
     \WedgeArrow{q0}{0}\WedgeArrow{q1}{0}\WedgeArrow{q3}{0}\WedgeArrow{q4}{0}\WedgeArrow{q2}{90}
     \begin{scope}[text=Zcol!85!black, font=\footnotesize]
        \node[below, yshift=-3pt] at (q1) {$Z_1$}; \node[right, xshift=2pt] at (q2) {$Z_2$}; \node[above, yshift=3pt] at (q4) {$Z_4^{-1}$};
     \end{scope}
  \end{scope}
  \begin{scope}[shift={(\x, \yBot)}]
     \FiveQPatchEdges
     \foreach \P/\txt in {q0/0, q3/3}{ \draw[line width=\qdLine,fill=white] (\P) circle (\qdR); \node[font=\numFont] at (\P) {\txt}; }
     \foreach \P/\txt in {q1/1, q2/2, q4/4}{ \draw[line width=\qdLine,fill=Zcol] (\P) circle (\qdR); \node[font=\numFont, text=white] at (\P) {\txt}; }
     \WedgeArrow{q0}{0}\WedgeArrow{q1}{0}\WedgeArrow{q3}{0}\WedgeArrow{q4}{0}\WedgeArrow{q2}{90}
     \begin{scope}[text=Zcol!85!black, font=\footnotesize]
        \node[below, yshift=-3pt] at (q1) {$Z_1$}; \node[right, xshift=2pt] at (q2) {$Z_2$}; \node[above, yshift=3pt] at (q4) {$Z_4^{-1}$};
     \end{scope}
  \end{scope}
  \draw[-{Stealth[length=3mm,width=2mm]},red!75,line width=1.6pt] (\x, 0.425) -- (\x, -0.1);

  % COL 3
  \def\x{1.6}
  \begin{scope}[shift={(\x, \yTop)}]
     \FiveQPatchEdges
     \foreach \P/\txt in {q3/3, q4/4}{ \draw[line width=\qdLine,fill=white] (\P) circle (\qdR); \node[font=\numFont] at (\P) {\txt}; }
     \foreach \P/\txt in {q0/0, q1/1, q2/2}{ \draw[line width=\qdLine,fill=Xcol] (\P) circle (\qdR); \node[font=\numFont, text=white] at (\P) {\txt}; }
     \WedgeArrow{q0}{0}\WedgeArrow{q1}{0}\WedgeArrow{q3}{0}\WedgeArrow{q4}{0}\WedgeArrow{q2}{90}
     \begin{scope}[text=Xcol!80!black, font=\footnotesize]
        \node[below, yshift=-3pt] at (q0) {$X_0$}; \node[below, yshift=-3pt] at (q1) {$X_1^{-1}$}; \node[right, xshift=2pt] at (q2) {$X_2$};
     \end{scope}
  \end{scope}
  \begin{scope}[shift={(\x, \yBot)}]
     \FiveQPatchEdges
     \foreach \P/\txt in {q3/3, q4/4}{ \draw[line width=\qdLine,fill=white] (\P) circle (\qdR); \node[font=\numFont] at (\P) {\txt}; }
     \foreach \P/\txt in {q0/0, q1/1, q2/2}{ \draw[line width=\qdLine,fill=Xcol] (\P) circle (\qdR); \node[font=\numFont, text=white] at (\P) {\txt}; }
     \WedgeArrow{q0}{0}\WedgeArrow{q1}{0}\WedgeArrow{q3}{0}\WedgeArrow{q4}{0}\WedgeArrow{q2}{90}
     \begin{scope}[text=Xcol!80!black, font=\footnotesize]
        \node[below, yshift=-3pt] at (q0) {$X_0$}; \node[below, yshift=-3pt] at (q1) {$X_1$}; \node[right, xshift=2pt] at (q2) {$X_2^{-1}$};
     \end{scope}
  \end{scope}
  \draw[-{Stealth[length=3mm,width=2mm]},red!75,line width=1.6pt] (\x, 0.425) -- (\x, -0.1);

  % COL 4
  \def\x{4.8}
  \begin{scope}[shift={(\x, \yTop)}]
     \FiveQPatchEdges
     \foreach \P/\txt in {q0/0, q1/1}{ \draw[line width=\qdLine,fill=white] (\P) circle (\qdR); \node[font=\numFont] at (\P) {\txt}; }
     \foreach \P/\txt in {q2/2, q3/3, q4/4}{ \draw[line width=\qdLine,fill=Xcol] (\P) circle (\qdR); \node[font=\numFont, text=white] at (\P) {\txt}; }
     \WedgeArrow{q0}{0}\WedgeArrow{q1}{0}\WedgeArrow{q3}{0}\WedgeArrow{q4}{0}\WedgeArrow{q2}{90}
     \begin{scope}[text=Xcol!80!black, font=\footnotesize]
        \node[right, xshift=2pt] at (q2) {$X_2^{-1}$}; \node[above, yshift=3pt] at (q3) {$X_3$}; \node[above, yshift=3pt] at (q4) {$X_4^{-1}$};
     \end{scope}
  \end{scope}
  \begin{scope}[shift={(\x, \yBot)}]
     \FiveQPatchEdges
     \foreach \P/\txt in {q0/0, q2/2}{ \draw[line width=\qdLine,fill=white] (\P) circle (\qdR); \node[font=\numFont] at (\P) {\txt}; }
     \foreach \P/\txt in {q1/1, q3/3, q4/4}{ \draw[line width=\qdLine,fill=Xcol] (\P) circle (\qdR); \node[font=\numFont, text=white] at (\P) {\txt}; }
     \WedgeArrow{q0}{0}\WedgeArrow{q1}{0}\WedgeArrow{q3}{0}\WedgeArrow{q4}{0}\WedgeArrow{q2}{90}
     \begin{scope}[text=Xcol!80!black, font=\footnotesize]
        \node[below, yshift=-3pt] at (q1) {$X_1^{-1}$}; \node[above, yshift=3pt] at (q3) {$X_3$}; \node[above, yshift=3pt] at (q4) {$X_4^{-1}$};
     \end{scope}
  \end{scope}
  \draw[-{Stealth[length=3mm,width=2mm]},red!75,line width=1.6pt] (\x, 0.425) -- (\x, -0.1);
\end{tikzpicture}
\caption{Stabilizer generators}
\end{subfigure}

\vspace{2.0ex}
% ===================== (e) Logical X =====================
\begin{subfigure}[c]{0.35\textwidth}
\centering
\begin{tikzpicture}[line cap=round,line join=round,thick,scale=1.0]
  \node[scale=1.0, anchor=east] at (-1.4, 0) {$\bar{X'}=\bar{X}\!:=$};
   
  % Added blue box
  \def\lw{1.5pt} \def\rc{8pt}
  \coordinate (LTL_X) at (-1.3, 1.4); \coordinate (LBR_X) at (1.4, -1.4);
  \coordinate (LTR_X) at ( 1.4, 1.4); \coordinate (LBL_X) at (-1.3, -1.4);
  \draw[Xcol, dash pattern=on 2pt off 5pt, line width=\lw, rounded corners=\rc] (LBL_X)--(LTL_X)--(LTR_X)--(LBR_X)--cycle;
   
  \FiveQPatchEdges
  \foreach \P/\txt in {q1/1, q2/2, q4/4}{ \draw[line width=\qdLine,fill=white] (\P) circle (\qdR); \node[font=\numFont] at (\P) {\txt}; }
  \foreach \P/\txt in {q0/0, q3/3}{ \draw[line width=\qdLine,fill=Xcol] (\P) circle (\qdR); \node[font=\numFont, text=white] at (\P) {\txt}; }
  \WedgeArrow{q0}{0}\WedgeArrow{q1}{0}\WedgeArrow{q3}{0}\WedgeArrow{q4}{0}\WedgeArrow{q2}{90}
  \begin{scope}[text=Xcol!80!black, font=\footnotesize]
     \node[below, yshift=-3pt] at (q0) {$X_0$};
     \node[above, yshift=3pt]  at (q3) {$X_3$};
  \end{scope}
\end{tikzpicture}
\caption{Logical $\bar X$}
\end{subfigure}
\quad
% ===================== (f) Logical Z =====================
\begin{subfigure}[c]{0.6\textwidth}
\centering
\begin{tikzpicture}[line cap=round,line join=round,thick,scale=1.0]
  % Define centers
  \coordinate (Left)  at (-2.2, 0);
  \coordinate (Right) at ( 2.2, 0);

  % Labels
  \node[scale=1.0, anchor=east] at ($(Left)+(-1.4, 0)$) {$\bar{Z'}\!:=$};
  \node[scale=1.0, anchor=east] at ($(Right)+(-1.4, 0)$) {$\bar{Z}\!:=$};
   
  % Added red boxes
  \def\lw{1.5pt} \def\rc{8pt}
  % Left Box
  \coordinate (LTL_Z) at ($(Left)+(-1.3, 1.4)$); \coordinate (LBR_Z) at ($(Left)+(1.4, -1.4)$);
  \coordinate (LTR_Z) at ($(Left)+(1.4, 1.4)$); \coordinate (LBL_Z) at ($(Left)+(-1.3, -1.4)$);
  \draw[Zcol, dash pattern=on 2pt off 5pt, line width=\lw, rounded corners=\rc] (LBL_Z)--(LTL_Z)--(LTR_Z)--(LBR_Z)--cycle;

  % Right Box
  \coordinate (RTL_Z) at ($(Right)+(-1.3, 1.4)$); \coordinate (RBR_Z) at ($(Right)+(1.4, -1.4)$);
  \coordinate (RTR_Z) at ($(Right)+(1.4, 1.4)$); \coordinate (RBL_Z) at ($(Right)+(-1.3, -1.4)$);
  \draw[Zcol, dash pattern=on 2pt off 5pt, line width=\lw, rounded corners=\rc] (RBL_Z)--(RTL_Z)--(RTR_Z)--(RBR_Z)--cycle;

  % Left Patch: Z0 Z1
  \begin{scope}[shift={(Left)}]
     \FiveQPatchEdges
     \foreach \P/\txt in {q2/2, q3/3, q4/4}{ \draw[line width=\qdLine,fill=white] (\P) circle (\qdR); \node[font=\numFont] at (\P) {\txt}; }
     \foreach \P/\txt in {q0/0, q1/1}{ \draw[line width=\qdLine,fill=Zcol] (\P) circle (\qdR); \node[font=\numFont, text=white] at (\P) {\txt}; }
     \WedgeArrow{q0}{0}\WedgeArrow{q1}{0}\WedgeArrow{q3}{0}\WedgeArrow{q4}{0}\WedgeArrow{q2}{90}
     \begin{scope}[text=Zcol!85!black, font=\footnotesize]
        \node[below, yshift=-3pt] at (q0) {$Z_0$};
        \node[below, yshift=-3pt] at (q1) {$Z_1$};
     \end{scope}
  \end{scope}

  % Right Patch: Z0 Z2
  \begin{scope}[shift={(Right)}]
     \FiveQPatchEdges
     \foreach \P/\txt in {q1/1, q3/3, q4/4}{ \draw[line width=\qdLine,fill=white] (\P) circle (\qdR); \node[font=\numFont] at (\P) {\txt}; }
     \foreach \P/\txt in {q0/0, q2/2}{ \draw[line width=\qdLine,fill=Zcol] (\P) circle (\qdR); \node[font=\numFont, text=white] at (\P) {\txt}; }
     \WedgeArrow{q0}{0}\WedgeArrow{q1}{0}\WedgeArrow{q3}{0}\WedgeArrow{q4}{0}\WedgeArrow{q2}{90}
     \begin{scope}[text=Zcol!85!black, font=\footnotesize]
        \node[below, yshift=-3pt] at (q0) {$Z_0$};
        \node[right, xshift=2pt]  at (q2) {$Z_2$};
     \end{scope}
  \end{scope}

  % Arrow in middle (adjusted for gap)
  \draw[-{Stealth[length=3mm,width=2mm]},red,line width=1.6pt] (-0.6, 0) -- (-0.1, 0);

\end{tikzpicture}
\caption{Logical $\bar Z$}
\end{subfigure}
% END OF PAGE 1
\end{figure*}

\begin{figure*}[tp]\ContinuedFloat
\centering
% ===== MAIN CAPTION HERE =====
% Caption shortened and scoped to the patch with q2 at the center (track hw, 2026-11-13). The former caption is in paper-rework/draft/.
\caption{The $[[5,1,2]]_3$ code on the patch with $q_2$ at the center. Primed symbols are the definitions of the folded-code construction~\cite{Moussa2016}, and unprimed symbols are ours, which exchange the roles of $q_1$ and $q_2$. Blue marks $X$-type and red $Z$-type operators; in (b) and (c), orange marks one-qutrit gates and purple the two-qutrit gate. The wedge on a qutrit points along its edge of the surface-code lattice (\appref{app:background}).
(a)~The patch of the folded-code construction, unfolded and folded, and our stabilizers.
(b)~Logical Hadamard gate: a layer of one-qutrit gates (Layer~1) and one SWAP gate (Layer~2), between the opposite corners $q_1$ and $q_3$ in $\bar H'$ and between the diagonal neighbors $q_2$ and $q_3$ in $\bar H$.
(c)~Logical phase gate, with its CZ gate on $q_1$ and $q_3$ in $\bar S'$ and on $q_2$ and $q_3$ in $\bar S$.
(d)~Stabilizer generators of the folded-code construction (top row) and ours (bottom row), $G_1$ to $G_4$ from left to right.
(e)~Logical $\bar X$, the same in both.
(f)~Logical $\bar Z' = Z_0 Z_1$ and $\bar Z = Z_0 Z_2$.}
\label{fig: 5-qutrit-code}

\end{figure*}

